% Generated from lessons/0001-open-closed-by-balls.html; edit the HTML source, then rerun the converter.
\section{用球邻域判断开闭}
\lessonmark{用球邻域判断开闭}

这节课只解决一个期末高频问题：给你 \(R^n\) 里的一个集合，怎样用“能不能放进一个小球”来判断它是开集、闭集、两者都是，还是两者都不是。

\subsection*{先抓住一句话}

在 \(R^n\) 中，开集的本质是：集合里的每一个点，都有一点“安全距离”。也就是从这个点出发，能放一个足够小的开球，并且这个球完全留在集合里。

\[
      E\text{ open}\quad\Longleftrightarrow\quad
      \forall x\in E,\ \exists r>0,\ B(x,r)\subset E
    \]

闭集不要硬背成“不开的反面”。闭集的常用判法是看补集是否开，或者看 \(E\) 中收敛序列的极限是否还留在 \(E\) 中。

\[
      E\text{ closed}
      \quad\Longleftrightarrow\quad
      R^n\setminus E\text{ open}
      \quad\Longleftrightarrow\quad
      (x_k\in E,\ x_k\to x \Rightarrow x\in E)
    \]

教材对应：《数学分析》第三版下册第十一章 §1“开集与闭集”。来源参考：\href{https://ocw.mit.edu/courses/18-100b-real-analysis-spring-2025/resources/ocw_18100b-lec13-2025apr03_mp4/}{MIT OCW Lecture 13} 和 \href{https://www.math.ucdavis.edu/~hunter/intro_analysis_pdf/ch13.pdf}{UC Davis Chapter 13}。

\subsection*{考试四步法}

\begin{conceptbox}
\textbf{1. 看边界}先找最容易靠近但可能不属于集合的点。
\end{conceptbox}

\begin{conceptbox}
\textbf{2. 判开}任取集合内一点，问能不能放一个完全留在集合内的小球。
\end{conceptbox}

\begin{conceptbox}
\textbf{3. 判闭}看边界点是否都被收进来，或改证补集是开集。
\end{conceptbox}

\begin{conceptbox}
\textbf{4. 写理由}开集写半径 \(r\)；闭集写补集开或序列极限。
\end{conceptbox}

期末答题时，最值号 < 往往对应“边界没收进来”，\(\le\) 往往对应“边界收进来”。但这只是提醒，不是证明。证明仍要回到球邻域、补集或序列。

\subsection*{例题：开圆盘}

设 \(A=\{(x,y)\in R^2:x^2+y^2<1\}\)。直觉上它是开集，因为任意内部点到圆周都有正距离。

证明时写成：

取 \(p\in A\)，令 \(\rho=\lVert p\rVert<1\)。取 \(r=(1-\rho)/2\)。若 \(\lVert q-p\rVert<r\)，则 \(\lVert q\rVert\le \lVert q-p\rVert+\lVert p\rVert<r+\rho<1\)，所以 \(q\in A\)。

它不是闭集，因为序列 \(p_k=(1-1/k,0)\) 都在 \(A\) 中，但 \(p_k\to(1,0)\)，而 \((1,0)\notin A\)。

\subsection*{马上练}

每题只选一种：开非闭、闭非开、开且闭、非开闭。选完立刻看反馈。

\begin{quickcheck}
\(A=\{(x,y):x^2+y^2<1\}\)

\tcblower
\textbf{答案：}开非闭。半径小于 1 的开圆盘是开集；边界圆周不在集合内，所以不是闭集。
\end{quickcheck}

\begin{quickcheck}
\(B=\{(x,y):x^2+y^2\le1\}\)

\tcblower
\textbf{答案：}闭非开。闭圆盘包含边界，所以是闭集；边界点周围任意小球都会跑到圆盘外，所以不是开集。
\end{quickcheck}

\begin{quickcheck}
\(C=\{(x,y):x>0,\ y\ge0\}\)

\tcblower
\textbf{答案：}非开闭。\(x=0\) 这条边界没有收进来，导致不闭；\(y=0\) 且 \(x>0\) 的点周围小球会跑到 \(y<0\)，导致不开。
\end{quickcheck}

\begin{quickcheck}
\(D=R^2\)

\tcblower
\textbf{答案：}开且闭。\(R^2\) 是全集，任意点周围的小球都在 \(R^2\) 内；它的补集是空集，也是开集，所以它也闭。
\end{quickcheck}

\subsection*{本节收束}

今天只要带走一个动作：判断开集时，不要先背结论，先任取集合内一点，然后试着造一个半径 \(r\)。判断闭集时，优先找“漏掉的极限点”，找不到再考虑补集是否开。

整节复习页：第十一章 §1：Euclid 空间上的基本定理。

下一课：聚点、闭包、边界。

教材路线：Euclid 空间上的极限和连续。

速查表：欧氏空间期末速查表。

推荐阅读：\href{https://ocw.mit.edu/courses/18-100b-real-analysis-spring-2025/mit18_100b_s25_lec_full.pdf}{MIT OCW 18.100B Real Analysis full lecture notes} 中关于 metric spaces、open and closed sets、compactness 的部分。

如果某一步卡住，直接把题目或你的证明发给我，我会按“概念点、判题依据、标准证明”帮你拆。
